\documentclass{article}
\usepackage[T1]{fontenc}
\usepackage{lmodern}
\usepackage{graphicx}
\usepackage{amsmath,amssymb}
\usepackage[a4paper,margin=2cm]{geometry}
\usepackage[absolute,overlay]{textpos}
\setlength{\TPHorizModule}{1mm}
\setlength{\TPVertModule}{1mm}
\usepackage[indonesian]{babel}
\usepackage{hyperref}
\setlength{\emergencystretch}{2em}
% unit: Halaman 17

\begin{document}

% === Halaman 17 (llm) ===

\section*{Serangan pada Diffie-Hellman}

\begin{itemize}
    \item Pertukaran kunci Diffie-Hellman hanya aman terhadap serangan pasif, misalnya penyadapan (\textit{eavesdropping}).
    \item Eva yang menyadap komunikasi antara Alice dan Bob hanya mengetahui nilai $g, p, A$ dan $B$ (semuanya publik). Dia tidak dapat mendeduksi $a$ dan $b$ (kunci privat Alice dan Bob) dari nilai-nilai publik tersebut. Untuk mendeduksi $a$ dan $b$, Eva harus dapat menghitung logaritma diskrit dari persamaan
    \[ A = g^a \pmod{p} \text{ atau } B = g^b \pmod{p} \]
    \item Namun, Diffie-Hellman tidak tahan terhadap serangan aktif seperti \textit{man-in-the-middle attack}.
    \item Misalkan Mallory mengintervensi komunikasi antara Alice dan Bob. Kepada Alice Malllory menyamar (\textit{impersonation}) sebagai Bob, dan kepada Bob Mallory menyamar sebagai Alice.
\end{itemize}

\vfill
\begin{center}
    Rinaldi M/II4020 Kriptografi
\end{center}

\end{document}
